\uuid{6GIR}
\titre{Changement de variables}
\auteur{hanani ; mzari}
\organisation{exo7}
\licence{CC BY-NC-SA 4.0}
\datecreate{2020-01-20}
\niveau{L1}
\module{Analyse}
\chapitre{Calcul d'intégrales}
\sousChapitre{Primitives diverses}
\difficulte{3}
\theme{Primitives, changement de variables}
\exo7qcmid{lille-2.primitives.029}

\contenu{
\begin{question}
\qclassification{127.11}
\qtags{motcle=changement de variables}
Le changement de variable $u=\sin x$ donne :
\begin{answers}
\good{$\displaystyle \int \frac{\cos x}{1+\sin ^2x}\mathrm{d}x=\int \frac{\mathrm{d}u}{1+u^2}$.}
\good{$\displaystyle \int \frac{\cos x}{1+\sin ^2x}\mathrm{d}x=\arctan (u)+k$, $k\in \Rr$.}
\bad{$\displaystyle \int \frac{\mathrm{d}x}{1+\sin ^2x}=\int \frac{\mathrm{d}u}{1+u^2}$.}
\bad{$\displaystyle \int \frac{\mathrm{d}x}{1+\sin ^2x}=\arctan (u)+k$, $k\in \Rr$.}
\end{answers}
\begin{explanations}
Avec $u=\sin x$, on aura $\displaystyle \mathrm{d}u=\cos x\mathrm{d}x$. D'où
$$\int \frac{\cos x}{1+\sin ^2x}\mathrm{d}x=\int \frac{\mathrm{d}u}{1+u^2}=\arctan (u)+k,\; k\in \Rr.$$
\end{explanations}
\end{question}
}
